\documentclass[11pt]{article}

\usepackage[margin=1in]{geometry}
\usepackage{amsmath}
\usepackage{amssymb}
\usepackage{booktabs}
\usepackage{graphicx}
\usepackage{hyperref}
\usepackage[numbers]{natbib}

\title{Inventory Control Policy Benchmarks for Selected Inventory Problems}
\author{Working manuscript}
\date{\today}

\begin{document}
\maketitle

\section{Problem focus}

This note records the current architecture study for the $L=4$, $p=4$ lost-sales family. The
central question is not whether one policy class is universally optimal, but which policy
parameterizations expose a replenishment rule that CMA-ES can search reliably.

The main numerical study below is still the lost-sales benchmark suite. We additionally record
three single-instance Rust-side benchmarks for perishable inventory, random-yield inventory, and
joint replenishment, using the same trained soft-tree policy family and the current
literature-verification rule: when public benchmark numbers are available, we report against them;
otherwise we report against our own exact solver and explicitly mark the instance as not verified
against a published table.

We study two closely related problems:
\begin{itemize}
    \item \textbf{vanilla lost sales}: no fixed setup cost,
    \item \textbf{fixed-cost lost sales}: a positive fixed ordering cost $K$ is charged whenever
    $q_t > 0$.
\end{itemize}

In both cases the controller observes the current inventory state and chooses an integer order
quantity $q_t \in \mathbb{Z}_{\ge 0}$. Demand is stochastic and unmet demand is lost. The
fixed-cost version adds the term $K\mathbf{1}\{q_t>0\}$ to the one-period cost, which makes the
order/no-order boundary much more important. Some policy families still introduce a finite
decoder-side cap $\bar q$, but that cap is now a policy design choice rather than an environment
restriction.

\section{Current benchmark instances and protocol}

The current single-instance study uses:
\begin{itemize}
    \item vanilla instance \texttt{vanilla\_l4\_p4\_poisson5},
    \item fixed-cost instance \texttt{lit\_pois\_mu5\_l4\_p4\_k5}.
\end{itemize}

Their shared core parameters are:
\[
L = 4, \qquad p = 4, \qquad h = 1, \qquad \mu = 5.
\]
The fixed-cost instance additionally uses $K=5$.

All current runs in this note use the same training budget:
\begin{itemize}
    \item $2000$ CMA-ES iterations,
    \item population size $64$,
    \item training horizon $2000$,
    \item evaluation horizon $10^6$,
    \item $10$ evaluation seeds,
    \item base random seed $42$ for training and heuristic search; reported mean and standard
    deviation values are evaluated over the consecutive seed set
    $\{42,43,\ldots,51\}$.
\end{itemize}

For the demand-robustness sweep we use four mean-preserving demand families:
\begin{itemize}
    \item Poisson demand with mean $5$,
    \item Geometric demand with mean $5$,
    \item positively correlated two-state Markov-modulated Poisson demand with
    $(\lambda_{\mathrm{low}}, \lambda_{\mathrm{high}})=(3,7)$ and $p_{00}=p_{11}=0.9$,
    \item negatively correlated two-state Markov-modulated Poisson demand with
    $(\lambda_{\mathrm{low}}, \lambda_{\mathrm{high}})=(3,7)$ and $p_{00}=p_{11}=0.1$.
\end{itemize}

For the two-state Markov-modulated Poisson cases, the latent regime $S_t \in \{\mathrm{low},
\mathrm{high}\}$ evolves as a two-state Markov chain with transition matrix
\[
P =
\begin{bmatrix}
p_{00} & 1-p_{00} \\
1-p_{11} & p_{11}
\end{bmatrix},
\]
and demand is then sampled as
\[
D_t \mid S_t=\mathrm{low} \sim \mathrm{Poisson}(\lambda_{\mathrm{low}}),
\qquad
D_t \mid S_t=\mathrm{high} \sim \mathrm{Poisson}(\lambda_{\mathrm{high}}).
\]
Because the reported cases use $p_{00}=p_{11}$, the stationary regime probabilities are
$\Pr(S_t=\mathrm{low})=\Pr(S_t=\mathrm{high})=1/2$, so both correlated variants preserve the same
unconditional mean demand:
\[
\mathbb{E}[D_t] = \tfrac{1}{2}\lambda_{\mathrm{low}} + \tfrac{1}{2}\lambda_{\mathrm{high}}
= \tfrac{1}{2}(3+7)=5.
\]
The sign and magnitude of temporal dependence come from the lag-one regime eigenvalue
\[
\rho_{\mathrm{regime}} = p_{00}+p_{11}-1.
\]
In the positive-correlation case this equals $0.8$, so the regime is persistent and the expected
run length in either state is $1/(1-0.9)=10$ periods. In the negative-correlation case it equals
$-0.8$, so the latent regime is strongly alternating. After Poisson observation noise is added,
the implemented lag-one demand autocorrelations become
\[
\mathrm{Corr}(D_t,D_{t+1})
= \frac{\pi_{\mathrm{low}}\pi_{\mathrm{high}}(\lambda_{\mathrm{high}}-\lambda_{\mathrm{low}})^2
(p_{00}+p_{11}-1)}{\mathrm{Var}(D_t)}
= \pm \frac{3.2}{9}
\approx \pm 0.356,
\]
with the plus sign for the persistent case and the minus sign for the alternating case. These are
the first-lag correlations checked explicitly in the Rust demand tests. More generally, the
implemented lag-$k$ demand autocovariance is proportional to $(p_{00}+p_{11}-1)^k$, so the
positive case produces geometrically decaying positive dependence and the negative case produces
geometrically decaying sign-alternating dependence. By contrast, the Poisson and Geometric i.i.d.
families have zero autocorrelation at every positive lag.

\section{Policy families and action equations}

The current study is intentionally narrow: the reported experiments use linear policies and soft
trees on the nonnegative scalar action space
\[
a_t \in \mathbb{Z}_{\ge 0}.
\]
Let $x_t \in \mathbb{R}^d$ denote the state feature vector. Some decoders additionally introduce a
finite policy-side cap $\bar q$ and therefore operate on the bounded control set
\[
\mathcal{A}_{\bar q} = \{0,1,\ldots,\bar q\},
\qquad
\mathcal{A}_{\bar q}^{+} = \{1,\ldots,\bar q\}.
\]
For uncapped nonnegative heads we write
\[
\mathcal{A}_{\infty} = \mathbb{Z}_{\ge 0}.
\]
The code implements a \emph{backbone} plus a \emph{decoder}:
\begin{itemize}
    \item a linear backbone emits logits by an affine map,
    \item a one-hidden-layer neural backbone emits logits through an MLP,
    \item a soft tree emits split probabilities and leaf-local quantity predictions.
\end{itemize}

For dense policies, let $u_\theta(x_t)$ denote the backbone output. In the linear case,
\[
u_\theta(x_t) = W x_t + b.
\]
In the one-hidden-layer neural case with width $H$ and activation $\phi$,
\[
u_\theta(x_t) = V \phi(Ux_t + c) + e.
\]
The decoder formulas below are shared by both dense backbones; only the way the latent outputs are
produced changes. For scalar controls we write $\Pi_{\mathcal{A}}(z)$ for nearest-integer
projection of a real-valued proposal $z$ onto the discrete action set $\mathcal{A}$.

We study the following action parameterizations:

\paragraph{Categorical quantity.}
The policy emits $\bar q+1$ logits and chooses
\[
a_t = \arg\max_{a \in \mathcal{A}_{\bar q}} u_a(x_t).
\]
This head is inherently cap-dependent.

\paragraph{Sigmoid direct quantity.}
The policy emits one scalar and maps it as
\[
a_t = \mathrm{round}\!\left(\bar q \, \sigma(u(x_t))\right).
\]
This is the code path named \texttt{sigmoid\_direct\_quantity}: a one-head bounded S-curve map.

\paragraph{Direct quantity.}
The policy emits one scalar and maps it as
\[
a_t = \Pi_{\mathcal{A}_{\infty}}\!\left(\mathrm{softplus}(u(x_t))\right).
\]
This is the code path named \texttt{direct\_quantity}: a one-head positive-part map with no
decoder-side upper bound.

\paragraph{Capped direct quantity.}
The policy emits one scalar and maps it as
\[
a_t = \Pi_{\mathcal{A}_{\bar q}}\!\left(\mathrm{softplus}(u(x_t))\right).
\]
This is the code path named \texttt{capped\_direct\_quantity}. It is the same softplus quantity
map as above, but with a decoder-side cap entering only through the final projection.

\paragraph{Soft-gated sigmoid direct quantity.}
The policy emits two latent outputs, a gate score $g(x_t)$ and a quantity score $q(x_t)$, then
sets
\[
a_t = \mathrm{round}\!\left(\sigma(g(x_t)) \, \sigma(q(x_t)) \, \bar q\right).
\]
This is the code path named \texttt{gated\_sigmoid\_direct\_quantity}.

\paragraph{Soft-gated direct quantity.}
The policy emits $g(x_t)$ and $q(x_t)$, then sets
\[
a_t = \Pi_{\mathcal{A}_{\bar q}}\!\left(\sigma(g(x_t)) \, \mathrm{softplus}(q(x_t))\right).
\]
This is the code path named \texttt{soft\_gated\_direct\_quantity}.

\paragraph{Hard-gated direct quantity.}
The policy emits $g(x_t)$ and $q(x_t)$, and uses
\[
a_t =
\begin{cases}
0, & \sigma(g(x_t)) < 0.5, \\
\Pi_{\mathcal{A}_{\bar q}^{+}}\!\left(\mathrm{softplus}(q(x_t))\right),
& \sigma(g(x_t)) \ge 0.5.
\end{cases}
\]
This is the code path named \texttt{hard\_gated\_direct\_quantity}.

\paragraph{Soft trees with linear leaves.}
A soft tree computes path probabilities $\pi_\ell(x_t)$ over leaves, and each leaf emits a local
linear quantity model
\[
q_\ell(x_t) = \mathrm{softplus}(\alpha_\ell^\top x_t + \beta_\ell).
\]
The final action is the rounded mixture
\[
a_t = \Pi_{\mathcal{A}_{\bar q}}\!\left(\sum_\ell \pi_\ell(x_t) q_\ell(x_t)\right).
\]
This is the active \texttt{soft\_tree\_depth1\_linear\_leaf} /
\texttt{soft\_tree\_depth2\_linear\_leaf} family. In code, the split network is always oblique in
the reported runs and the scalar leaf is exactly a softplus-transformed affine quantity rule.

The policy family that is \emph{not} part of the active benchmark surface is the raw unbounded
direct head
\[
a_t = \Pi_{\mathcal{A}_{\bar q}}\!\left(u(x_t)\right),
\]
which repeatedly converged to weak basins in earlier runs and is now treated as a negative ablation
rather than a serious candidate.

\subsection{Parameter counts}

The policy families discussed in this note fall into three groups: linear dense policies,
one-hidden-layer neural dense policies, and soft trees. For the active $L=4$ pipeline state we
have $d=4$ and therefore $d+1=5$. The bounded heads use $\bar q=20$ in the canonical vanilla
instance and $\bar q=50$ in the canonical fixed-cost instance. The current NN registry entries use
a single hidden layer of width $H=50$ with SELU activation. We exclude the ordinal decoders from
the active rerun surface below because their output dimension, and hence their parameter count,
still scales directly with $\bar q$.

\begin{table}[t]
\centering
\small
\caption{Trainable parameter counts for the linear scalar-quantity policy families. The last
column shows the concrete $L=4$ counts, reported as vanilla / fixed-cost when the decoder-side
cap differs across the two canonical instances.}
\label{tab:linear-parameter-counts}
\begin{tabular}{|l|c|c|}
\hline
Policy family & Trainable parameters & Count at $L=4$ \\
\hline
Linear categorical quantity & $(d+1)(\bar q+1)$ & $105 / 255$ \\
Linear sigmoid direct quantity & $d+1$ & $5$ \\
Linear direct quantity & $d+1$ & $5$ \\
Linear capped direct quantity & $d+1$ & $5$ \\
Linear gated sigmoid direct quantity & $2(d+1)$ & $10$ \\
Linear soft-gated direct quantity & $2(d+1)$ & $10$ \\
Linear hard-gated direct quantity & $2(d+1)$ & $10$ \\
Soft tree, depth-1 linear leaf & $3(d+1)$ & $15$ \\
Soft tree, depth-2 linear leaf & $7(d+1)$ & $35$ \\
\hline
\end{tabular}
\end{table}

\begin{table}[t]
\centering
\small
\caption{Trainable parameter counts for the one-hidden-layer NN scalar-quantity policy families.
These counts use the current repo default $H=50$ for the \texttt{nn\_*} policies. The last column
again reports vanilla / fixed-cost when the decoder-side cap differs across the two canonical
instances.}
\label{tab:nn-parameter-counts}
\begin{tabular}{|l|c|c|}
\hline
Policy family & Trainable parameters & Count at $L=4$ \\
\hline
NN categorical quantity & $H(d+1) + (H+1)(\bar q+1)$ & $1321 / 2851$ \\
NN sigmoid direct quantity & $H(d+1) + (H+1)$ & $301$ \\
NN direct quantity & $H(d+1) + (H+1)$ & $301$ \\
NN capped direct quantity & $H(d+1) + (H+1)$ & $301$ \\
NN gated sigmoid direct quantity & $H(d+1) + 2(H+1)$ & $352$ \\
NN soft-gated direct quantity & $H(d+1) + 2(H+1)$ & $352$ \\
NN hard-gated direct quantity & $H(d+1) + 2(H+1)$ & $352$ \\
\hline
\end{tabular}
\end{table}

\begin{table}[t]
\centering
\small
\caption{Trainable parameter counts for the soft-tree policy families.}
\label{tab:tree-parameter-counts}
\begin{tabular}{|l|c|c|}
\hline
Policy family & Trainable parameters & Count at $L=4$ \\
\hline
Soft tree, depth-1 linear leaf & $3(d+1)$ & $15$ \\
Soft tree, depth-2 linear leaf & $7(d+1)$ & $35$ \\
\hline
\end{tabular}
\end{table}

\section{Numerical experiments}

\subsection{Vanilla lost sales}

Table~\ref{tab:vanilla-l4-demand-table} reports the full $L=4$, $p=4$ vanilla comparison across
all four demand families on the active linear/tree benchmark surface. All heuristic baselines are
shown explicitly rather than being collapsed into a single best-heuristic row.

\begin{table}[t]
\centering
\small
\caption{Vanilla lost-sales mean costs for the single-instance $L=4$, $p=4$ study across all
four demand families under the current $2000$-iteration / population-$64$ CMA-ES protocol. Bold
entries mark the best value in each demand column.}
\label{tab:vanilla-l4-demand-table}
\begin{tabular}{|l||c|c||c|c|}
\hline
 & \multicolumn{2}{c||}{i.i.d. demand} & \multicolumn{2}{c|}{Correlated demand} \\
\hline
Policy / heuristic & Poisson & Geometric & MMPP2+ & MMPP2- \\
\hline
myopic1 & 5.0627 & 11.3176 & 9.8925 & 5.7245 \\
myopic2 & 4.8168 & 10.8025 & 9.5016 & 5.4150 \\
SVBS & 5.8344 & 13.8596 & 12.2518 & 6.8244 \\
\hline
Linear categorical quantity & 4.8752 & 10.8041 & 8.6411 & 5.7431 \\
Linear sigmoid direct quantity & 4.7557 & 10.6606 & 8.6336 & 5.3673 \\
Linear capped direct quantity & 4.7467 & 10.6426 & 8.5899 & 5.3654 \\
Linear gated sigmoid direct quantity & \textbf{4.7451} & 10.6469 & 8.6206 & 5.3680 \\
Linear soft-gated direct quantity & 4.7481 & 10.6416 & 8.5819 & 5.3732 \\
Linear hard-gated direct quantity & 4.7480 & 10.6542 & 8.5799 & 5.3685 \\
Soft tree, depth-1 linear leaf & 4.7481 & \textbf{10.6391} & 8.5893 & --- \\
Soft tree, depth-2 linear leaf & 4.7518 & 10.6453 & 11.9015 & --- \\
\hline
\end{tabular}
\end{table}

The vanilla side is comparatively forgiving:
\begin{itemize}
    \item every non-categorical policy beats the best heuristic in the Poisson and Geometric
    cases,
    \item the correlated-demand winners remain compact dense direct heads rather than the trees,
    \item the one-head sigmoid direct policy is competitive at population $64$,
    \item the currently stored negative-MMPP2 vanilla bundle does not contain the two tree rows,
    so those cells are marked unavailable.
\end{itemize}

So the full vanilla numerical experiment still says the same thing: once we move away from the
flat categorical head, most direct and tree parameterizations land in essentially the same good
basin. Search geometry matters here, but only weakly compared with the fixed-cost case below.

\subsection{Fixed-cost lost sales}

Table~\ref{tab:fixed-cost-l4-demand-table} reports the corresponding fixed-cost $L=4$, $p=4$,
$K=5$ comparison on the same active linear/tree benchmark surface.

\begin{table}[t]
\centering
\small
\caption{Fixed-cost lost-sales mean costs for the single-instance $L=4$, $p=4$, $K=5$ study
across all four demand families under the current $2000$-iteration / population-$64$ CMA-ES
protocol. Bold entries mark the best value in each demand column.}
\label{tab:fixed-cost-l4-demand-table}
\begin{tabular}{|l||c|c||c|c|}
\hline
 & \multicolumn{2}{c||}{i.i.d. demand} & \multicolumn{2}{c|}{Correlated demand} \\
\hline
Policy / heuristic & Poisson & Geometric & MMPP2+ & MMPP2- \\
\hline
$(s,S)$ & 9.3677 & 14.1229 & 12.1579 & 9.8755 \\
$(s,nQ)$ & 9.2094 & 14.0669 & 12.2139 & 9.6925 \\
modified $(s,S,q)$ & 9.2003 & 14.0023 & 12.0935 & 9.6501 \\
\hline
Linear categorical quantity & 10.2754 & 15.9951 & 14.4001 & 10.7435 \\
Linear sigmoid direct quantity & 8.7728 & 15.6229 & 12.3011 & 9.0772 \\
Linear capped direct quantity & 9.7482 & 13.8152 & 12.4235 & 10.3640 \\
Linear gated sigmoid direct quantity & 9.7474 & 13.7711 & 13.2140 & 10.3766 \\
Linear soft-gated direct quantity & 8.7715 & 14.2049 & 12.3907 & \textbf{9.0656} \\
Linear hard-gated direct quantity & 9.7478 & \textbf{13.6276} & \textbf{11.8787} & 10.3620 \\
Soft tree, depth-1 linear leaf & 8.7708 & 13.8218 & 12.3165 & 10.3697 \\
Soft tree, depth-2 linear leaf & \textbf{8.7700} & 13.7412 & 12.3836 & 10.3675 \\
\hline
\end{tabular}
\end{table}

This is still the sharper architecture test:
\begin{itemize}
    \item there is no single winner across all fixed-cost demand families,
    \item flat categorical quantity remains weak in every case,
    \item positively correlated and overdispersed fixed-cost demand favor the hard-gated direct
    head,
    \item negatively correlated fixed-cost demand favors the soft-gated direct head,
    \item Poisson fixed-cost demand still slightly favors the depth-$2$ tree.
\end{itemize}

This is where the full numerical experiment becomes much more informative than a single summary
table. In fixed-cost lost sales, changing the action head changes which basin CMA-ES reaches. The
problem is not just representational capacity: the main empirical effect is that different heads
expose different order/no-order geometry to the search procedure.

\section{Additional single-instance benchmarks}

This section records the first single-instance learned-policy benchmarks for the other Rust-side
problems we have already validated. The goal here is narrower than in the lost-sales study above:
for each problem we state the control model, the learned policy class, and the benchmark outcome on
the current canonical instance.

\subsection{Perishable inventory}

\paragraph{Formulation.}
We use the age-structured single-item perishable inventory model with stochastic demand, positive
lead time, and issuing discipline fixed to FIFO for the benchmark instance. At period $t$, the
state contains the pipeline orders and the on-hand inventory split by remaining lifetime. After the
controller chooses an order quantity $q_t \in \mathbb{Z}_{\ge 0}$, demand is realized, units age,
and expired stock is discarded. The one-period cost is
\[
c q_t + h H_t + p B_t + w W_t,
\]
where $H_t$ is end-of-period holding inventory, $B_t$ is unmet demand, and $W_t$ is waste. The
reported benchmark numbers are mean discounted return, so higher is better.

\paragraph{Policy.}
The learned policy is a depth-$2$ oblique soft decision tree with linear leaves. Each leaf outputs
an affine scalar score in the normalized state features, and the scalar action head uses the
uncapped nonnegative linear decoder already adopted on the Rust side for scalar-action problems.
For this benchmark the trained policy is \texttt{d=2, split=oblique, leaf=linear}.

\paragraph{Results.}
The benchmark instance is \texttt{de\_moor2022\_m4\_exp6\_l2\_cp7\_fifo} with shelf life $4$, lead
time $2$, demand mean $4$, coefficient of variation $0.5$, shortage cost $5$, holding cost $1$,
waste cost $7$, and procurement cost $3$. This problem family is literature-verified: on two
smaller De Moor verification instances, our exact value-iteration return and best base-stock level
match the published tables exactly. On the medium benchmark instance, the trained tree improves on
the two heuristic baselines but still trails the published exact value-iteration benchmark.

\begin{table}[t]
\centering
\caption{Single-instance perishable-inventory benchmark on
\texttt{de\_moor2022\_m4\_exp6\_l2\_cp7\_fifo}. Higher mean discounted return is better.}
\label{tab:perishable-single-instance}
\begin{tabular}{lrr}
\toprule
Policy & Mean discounted return & Gap vs.\ soft tree \\
\midrule
published value iteration & -1432.000 & 14.675 \\
published base stock & -1453.000 & -6.325 \\
repo base stock & -1454.363 & -7.688 \\
repo BSP-low-EW & -1454.363 & -7.688 \\
soft tree, depth-2 linear leaf & \textbf{-1446.675} & 0.000 \\
\bottomrule
\end{tabular}
\end{table}

\subsection{Random-yield inventory}

\paragraph{Formulation.}
We use the single-item random-yield inventory problem with backlogging, positive lead time, and
all-or-nothing delivery success. The state consists of the current inventory level and the pipeline
orders. When the order placed $L$ periods earlier reaches the head of the pipeline, it arrives in
full with probability $\rho$ and otherwise fails completely. After arrival and demand realization,
the one-period cost is
\[
c q_t + h I_t^+ + p (-I_t)^+,
\]
where $I_t$ is the end-of-period inventory position. The reported benchmark numbers are mean
discounted cost, so lower is better.

\paragraph{Policy.}
The learned policy is a depth-$1$ oblique soft decision tree with linear leaves, again using the
uncapped scalar linear decoder on the Rust side. The trained benchmark policy is
\texttt{d=1, split=oblique, leaf=linear}. It is trained with seed-batched rollouts and then
re-evaluated on a disjoint holdout seed set.

\paragraph{Results.}
The benchmark instance is \texttt{yan2026\_style\_lt2\_p075\_discounted} with horizon $12$, lead
time $2$, demand mean $4$, success probability $0.75$, holding cost $1$, shortage cost $9$,
procurement cost $1$, discount factor $0.99$, initial inventory $6$, and initial pipeline
$(4,3)$. This is a repo-native instance shaped after the small-lead-time discounted regime
emphasized by Yan et al., not a verbatim published table row. We therefore do \emph{not} claim
literature verification for this benchmark instance; instead, correctness is anchored by a smaller
exact-DP verifier in the same model family, where the repo exact solution and heuristic costs are
frozen from the solver and match the stored assertions exactly.

\begin{table}[t]
\centering
\caption{Single-instance random-yield benchmark on
\texttt{yan2026\_style\_lt2\_p075\_discounted}. Lower mean discounted cost is better.}
\label{tab:random-yield-single-instance}
\begin{tabular}{lrr}
\toprule
Policy & Mean discounted cost & Gap vs.\ soft tree cost \\
\midrule
linear inflation & 203.762 & 12.022 \\
weighted newsvendor & 222.881 & 31.142 \\
soft tree, depth-1 linear leaf & \textbf{191.740} & 0.000 \\
\bottomrule
\end{tabular}
\end{table}

\subsection{Joint replenishment}

\paragraph{Formulation.}
We use the finite-horizon discounted joint replenishment problem with two items, shared truck
capacity, item-specific minor ordering costs, and holding and shortage penalties. At each period
$t$ the controller selects an order vector $q_t \in \mathbb{Z}_{\ge 0}^n$. If the total ordered
quantity exceeds one truckload, multiple major orders are charged, so the ordering term is
\[
K \left\lceil \frac{\sum_{i=1}^n q_{i,t}}{C} \right\rceil
+ \sum_{i=1}^n k_i \mathbf{1}\{q_{i,t} > 0\}.
\]
After stochastic demand is realized independently for each item, holding and shortage costs are
charged on the ending inventories. The reported benchmark numbers are mean discounted cost, so
lower is better.

\paragraph{Policy.}
The learned policy is a depth-$2$ oblique soft decision tree with linear leaves. Each leaf outputs
an affine vector in the normalized state features, and the resulting action is projected back into
the feasible order box because this problem has vector-valued bounded actions. For this benchmark
the trained policy is \texttt{d=2, split=oblique, leaf=linear}.

\paragraph{Results.}
The current benchmark instance is the repo exact verifier built from the small-scale setting family
used by Vanvuchelen et al.: horizon $4$, discount factor $0.99$, truck capacity $6$, maximum order
quantities $(12,12)$, initial inventory $(2,0)$, major order cost $75$, minor order costs
$(10,10)$, holding costs $(1,1)$, shortage costs $(19,19)$, and independent uniform item demands
on $\{0,\dots,5\}$ and $\{0,\dots,3\}$. The public paper reports the setting family but not a
public exact benchmark row for repo assertions, so this instance is \emph{not} literature-verified;
the optimal DP cost and heuristic benchmarks are instead frozen from our exact solver. On this
instance the trained tree remains above the exact optimum but substantially outperforms the MOQ and
DYN-OUT heuristics.

\begin{table}[t]
\centering
\caption{Single-instance joint-replenishment benchmark on the repo exact verification instance.
Lower mean discounted cost is better.}
\label{tab:joint-single-instance}
\begin{tabular}{lrr}
\toprule
Policy & Mean discounted cost & Gap vs.\ soft tree cost \\
\midrule
optimal DP & 266.386 & -38.122 \\
minimum order quantity & 386.101 & 81.593 \\
dynamic order-up-to & 383.960 & 79.452 \\
soft tree, depth-2 linear leaf & \textbf{304.508} & 0.000 \\
\bottomrule
\end{tabular}
\end{table}

\section{Interpretation}

The original local hypothesis was that \emph{factorization} or \emph{gating} might be the main
winning ingredient. The numerical evidence is now weaker and more precise than that.

So the best current reading is:
\begin{quote}
Good policies are the ones whose parameterization makes the relevant replenishment map easy enough
for CMA-ES to find under the available budget.
\end{quote}

That map need not always be represented in the same way. Vanilla lost sales is comparatively
forgiving, so many non-categorical direct and tree heads realize essentially the same solution.
Fixed-cost lost sales is not: there the action-head geometry and the demand family much more
strongly determine which replenishment boundary CMA-ES can actually find.

\subsection{Dual sourcing}

The current dual-sourcing study points in the same overall direction, but even more sharply:
policy \emph{design} matters more than raw flexibility.

For the hard $l_r=3$ rows in the Gijs Figure~9 family, we compared several structured learned
policies that all use the same reduced pipeline state and the same CMA-ES screening budget, but
different action parameterizations. The best current family is not the widest tree. It is a
small axis-aligned soft tree with constant leaves over a capped-dual-index-delta control grid.

The relevant learned control family is
\[
(s_e,\Delta_r,\bar c_r),
\qquad
s_r = s_e + \Delta_r,
\]
with the regular cap restricted to the small discrete set
\[
\bar c_r \in \{1,2,3,4,6,8,12\}.
\]
This keeps the learned policy close to the heuristic geometry while still allowing the targets to
depend on state.

At $l_r=3$ the reduced-state input dimension is $d=3$ and the control dimension is also $3$, so
the tree parameter counts are tiny:
\begin{itemize}
    \item depth-1 axis-aligned constant tree: $10$ parameters,
    \item depth-2 axis-aligned constant tree: $24$ parameters,
    \item depth-3 axis-aligned constant tree: $52$ parameters,
    \item depth-2 oblique linear tree on the same control family: $60$ parameters.
\end{itemize}

The performance differences are large. On the same repo evaluation protocol, the $60$-parameter
oblique linear small-cap tree yields gaps of about $0.40\%$ on \texttt{dual\_l3\_ce105} and
$1.59\%$ on \texttt{dual\_l3\_ce110} relative to the best heuristic. By contrast, the
$24$-parameter axis-aligned constant version reduces those to about $0.026\%$ and $0.209\%$.
Moreover, a $10$-parameter depth-1 axis-aligned constant tree already beats the best heuristic on
\texttt{dual\_l3\_ce105} under the screening protocol.

So the dual-sourcing evidence supports a stronger version of the same design lesson:
\begin{quote}
Small learned policies can match or outperform strong benchmark heuristics, but only when their
action parameterization respects the problem structure. Larger, more flexible trees can perform
materially worse when that structure is too loose.
\end{quote}

\section{Published verification anchor}

Before relying on the fixed-cost simulator and heuristic stack for learned-policy comparisons, we
validate the implementation against the published example used in
\citet{Bijvank2015ParametricPolicies}. Their Table~1 reports a small Poisson-demand instance with:
\[
L=2,\qquad p=14,\qquad K=5,\qquad h=1,\qquad \mu=5.
\]

In the repo, this validation anchor is:
\begin{itemize}
    \item \texttt{bijvank2015\_table1\_l2\_p14\_k5}
\end{itemize}
and is encoded in
\texttt{invman/problems/lost\_sales\_fixed\_order\_cost/reference\_instances.py}.

\begin{table}[t]
\centering
\caption{Published fixed-cost lost-sales validation example from
\citet{Bijvank2015ParametricPolicies} and our reproduced long-run costs for the reported heuristic
parameters.}
\label{tab:fixed-cost-validation}
\begin{tabular}{llll}
\toprule
Policy & Reported parameters & Published cost & Our reproduced cost \\
\midrule
optimal & --- & 11.46 & --- \\
$(s,S)$ & $(17,23)$ & 11.62 & 11.6217 \\
$(s,nQ)$ & $(17,7)$ & 11.56 & 11.5589 \\
modified $(s,S,q)$ & $(17,23,7)$ & 11.50 & 11.5012 \\
\bottomrule
\end{tabular}
\end{table}

The reproduced costs match the published values to within simulation precision. This is the
literature-backed correctness anchor for the fixed-cost environment and heuristic implementation.

It is important to distinguish this validation role from the architecture study above:
\begin{itemize}
    \item the validation instance is used to verify that the simulator and heuristic policies match
    the published benchmark family,
    \item the canonical $L=4$, $p=4$ studies are where we compare modern learned-policy
    parameterizations.
\end{itemize}

\bibliographystyle{plainnat}
\bibliography{references}

\end{document}
